\chapter{X-ray diffraction (XRD)}
\label{ch:xrd}

\paragraph{In this chapter.} Lattice planes and Miller indices, Bragg's law, d-spacings, systematic absences (how to tell
fcc from bcc), the full procedure to index a pattern and get the lattice parameter, peak broadening and the Scherrer
equation. XRD is the most frequent calculation in the exercises.

\section{Lattice planes and Miller indices}
\slide{169--171}: a crystal lattice contains many families of parallel planes, labelled by Miller indices $(hkl)$.
Each family gives one diffraction peak; its \textbf{d-spacing} sets the peak angle.
\begin{itemize}
  \item $(100)$, $(010)$, $(001)$: cube faces. $(110)$: diagonal plane. $(111)$: plane cutting the three axes at 1.
  \item $(200)$ has half the spacing of $(100)$; $(300)$ one third \slide{170}.
  \item $\{hkl\}$ = family of equivalent planes (e.g.\ $\{100\}$ = all six cube faces).
\end{itemize}

\begin{keybox}[: d-spacing for orthogonal systems \slide{172--173}]
\begin{equation}
\text{Cubic:}\quad \frac{1}{d_{hkl}^2} = \frac{h^2+k^2+l^2}{a^2}
\qquad\Longleftrightarrow\qquad \boxed{d_{hkl} = \frac{a}{\sqrt{h^2+k^2+l^2}}}
\label{eq:dcubic}
\end{equation}
Tetragonal: $\dfrac{1}{d^2} = \dfrac{h^2+k^2}{a^2} + \dfrac{l^2}{c^2}$; \quad
Orthorhombic: $\dfrac{1}{d^2} = \dfrac{h^2}{a^2}+\dfrac{k^2}{b^2}+\dfrac{l^2}{c^2}$; \quad
Hexagonal: $\dfrac{1}{d^2} = \dfrac{4}{3}\dfrac{h^2+hk+k^2}{a^2}+\dfrac{l^2}{c^2}$.
\end{keybox}

\section{Bragg's law}
\begin{keybox}[: Bragg equation \slide{174}]
\begin{equation}
\boxed{2\,d\,\sin\theta = n\,\lambda}
\label{eq:bragg}
\end{equation}
Rays reflected by consecutive planes have a path difference $2d\sin\theta$; constructive interference when it is an
integer number of wavelengths. Lab source: \textbf{Cu K$\alpha$, $\lambda = \SI{1.54}{\angstrom}$} \slide{168}.
\end{keybox}
\begin{pitfall}
The diffractogram's x-axis is $\twotheta$. Divide by 2 before taking $\sin\theta$. Using $\sin(2\theta)$ is the most
common error.
\end{pitfall}
Combining \cref{eq:dcubic,eq:bragg} ($n=1$):
\begin{equation}
\sin^2\theta = \frac{\lambda^2}{4a^2}\,(h^2+k^2+l^2).
\label{eq:sin2}
\end{equation}
Since $\lambda^2/4a^2$ is the same for all peaks, \textbf{$\sin^2\theta$ of the peaks are in the ratio of the integers
$N = h^2+k^2+l^2$}. The \emph{lowest-angle} peaks come from the \emph{largest} d-spacings (smallest $N$) \slide{177}.

\textbf{Monocrystal vs polycrystal} \slide{175--176}: a single crystal gives spots; a powder (random crystallites) gives
Debye--Scherrer rings, recorded by the diffractometer as peaks vs $\twotheta$.

\section{Systematic absences: which peaks appear}
\begin{keybox}[: selection rules for cubic lattices \slide{179--182}]
\begin{center}
\begin{tabular}{@{}lll@{}}
\toprule
Lattice & Reflection observed if\ldots & First lines \\ \midrule
Simple cubic (P) & all $hkl$ & 100, 110, 111, 200, 210, 211, \ldots \\
Body-centred (I, bcc) & $h+k+l$ even & 110, 200, 211, 220, 310, 222, \ldots \\
Face-centred (F, fcc) & $h,k,l$ all odd or all even & 111, 200, 220, 311, 222, 400, \ldots \\
\bottomrule
\end{tabular}
\end{center}
Why: in bcc the atom at the body centre lies on a plane halfway between two (100) planes and scatters exactly out of
phase, so (100) vanishes but (200) is observed \slide{179}. In fcc (NaCl) the (110) planes have an identical plane
halfway between them, so (110) is absent, while (111) is observed \slide{180}.
\end{keybox}

\begin{keybox}[: the fingerprint sequences of $N = h^2+k^2+l^2$]
\begin{center}
\small
\begin{tabular}{@{}lll>{\raggedright\arraybackslash}p{2.8cm}@{}}
\toprule
Lattice & $N$ sequence & ratios $N/N_1$ & looks like \\ \midrule
SC  & 1, 2, 3, 4, 5, 6, 8, 9 & 1:2:3:4:5:6:8 & integers, 7 missing \\
BCC & 2, 4, 6, 8, 10, 12, 14 & 1:2:3:4:5:6:7 & all integers \\
FCC & 3, 4, 8, 11, 12, 16, 19 & 1:1.33:2.67:3.67:4:5.33 & \textbf{pair, gap, pair} \\
\bottomrule
\end{tabular}
\end{center}
Visual pattern: fcc peaks come in \textbf{pairs} (111/200, 311/222, 331/420) with the first of the pair stronger;
bcc peaks are roughly evenly spaced \slide{183}.
\end{keybox}
\begin{tipbox}
SC and BCC both give integer ratios for the first six lines. Look at the 7th line (SC has no $N=7$, BCC gives 7 $\to$
$N=14$), or use physical knowledge: metals are fcc or bcc, almost never sc.
\end{tipbox}
Structures with a basis (e.g.\ ZnS zincblende, NaCl rocksalt, CeO$_2$ fluorite, TiN rocksalt) are on an \textbf{fcc}
lattice and follow the fcc rule. Some fcc peaks are weak or missing for compounds (e.g.\ ZnS (200) weak) \slide{184}.

\section{Procedure: indexing a cubic pattern}
\label{sec:xrd-indexing}
\begin{recipe}{indexing and lattice parameter}
\begin{enumerate}
  \item Read all peak positions $\twotheta_i$. Compute $\theta_i$ and $\sin^2\theta_i$.
  \item Divide by the first: $r_i = \sin^2\theta_i/\sin^2\theta_1$.
        \begin{itemize}
          \item $r = 1, 2, 3, 4, 5, 6 \Rightarrow$ BCC ($N = 2r$).
          \item $r = 1, 1.33, 2.67, 3.67, 4 \Rightarrow$ FCC ($N = 3r$).
        \end{itemize}
  \item Assign $hkl$ from $N$ (table above).
  \item For each peak: $d_i = \lambda/(2\sin\theta_i)$ and $a_i = d_i\sqrt{N_i}$.
  \item Average the $a_i$ (high-angle peaks are more precise).
\end{enumerate}
Sanity check: metals have $a\approx\SIrange{3}{4.2}{\angstrom}$; fluorite oxides (CeO$_2$) $\approx$ \SI{5.4}{\angstrom}.
\end{recipe}
Worked example: a cubic lattice with $a = \SI{4}{\angstrom}$, $\lambda = \SI{1.54}{\angstrom}$ gives (100) $d=4.00$, $\twotheta=22.2^\circ$;
(110) 2.83, $31.6^\circ$; (111) 2.31, $39.0^\circ$; (200) 2.00, $45.2^\circ$; (210) 1.79, $51.0^\circ$ \slide{178}.

\subsection*{Identifying a phase / distinguishing two compounds}
Compute the expected $\twotheta$ of the first allowed reflections for each candidate (from its $a$ and lattice type)
and compare with the measured pattern (Exercise~4). Two phases in a mixture give \textbf{two sets of peaks}; a
\textbf{solid solution} gives \textbf{one set} at intermediate positions (Exercise~6, \cref{sec:vegard}).

\section{Peak broadening and the Scherrer equation}
Real crystals are finite, so destructive interference is incomplete over a range of angles and peaks broaden
\slide{185}. Contributions to the observed width: \textbf{instrumental profile}, \textbf{crystallite size},
\textbf{microstrain}. Smaller crystallites give broader peaks (7 $\to$ 2 nm, \slide{186}).
\begin{keybox}[: Scherrer equation \slide{186}]
\begin{equation}
\boxed{L = \frac{K\,\lambda}{B_{(2\theta)}\,\cos\theta_\mathrm{max}}}, \qquad K = 0.9
\label{eq:scherrer}
\end{equation}
$L$: mean crystallite size; $K$: shape factor; $\theta_\mathrm{max}$: Bragg angle of the peak (half of its $\twotheta$);
$B$: FWHM of the peak \textbf{in radians}, after subtracting the instrumental broadening.
Class note: ``As particle size increases, peaks get narrower.''
\end{keybox}
\begin{pitfall}
\begin{enumerate}
  \item Convert $B$ to \textbf{radians}: $B[\mathrm{rad}] = B[^\circ]\cdot\pi/180$.
  \item Use $\cos\theta$ with $\theta = (\twotheta)/2$, not $\cos 2\theta$.
  \item Subtract the instrumental broadening: the course says ``after subtracting'', i.e.\ $B = B_\mathrm{obs}-B_\mathrm{inst}$.
        Some books use $B^2 = B_\mathrm{obs}^2 - B_\mathrm{inst}^2$\extra. Say which one you use.
  \item The result is in the units of $\lambda$ (usually \AA; divide by 10 for nm).
\end{enumerate}
\end{pitfall}
\begin{keybox}[: crystallite size $\neq$ particle size \slide{187}]
A particle may be made of several crystallites. Crystallite size often matches grain size, but not always.
XRD (Scherrer) gives the \emph{crystallite} size; TEM gives the \emph{particle} size. If $L_\mathrm{XRD} < d_\mathrm{TEM}$,
the particles are polycrystalline.
\end{keybox}

\section{Vegard's law (mixtures vs solid solutions)\extra}
\label{sec:vegard}
For a substitutional solid solution A$_{1-x}$B$_x$ with the same structure, the lattice parameter varies approximately
linearly between the end members:
\begin{equation}
a_{\mathrm{A}_{1-x}\mathrm{B}_x} \approx (1-x)\,a_\mathrm{A} + x\,a_\mathrm{B}
\qquad\Rightarrow\qquad x \approx \frac{a_\mathrm{A} - a_\mathrm{measured}}{a_\mathrm{A}-a_\mathrm{B}}.
\end{equation}
Not in the slides, but it is the standard way to answer ``mixture or solid solution?'' quantitatively (Exercise~6).
\begin{pitfall}
If the two lattice parameters are almost equal (Au \SI{4.065}{\angstrom} vs Ag \SI{4.079}{\angstrom}), XRD \textbf{cannot}
distinguish alloy from mixture. Use another technique (UV-vis, XPS, EELS, Exercise~16).
\end{pitfall}

\section{Advantages and limitations}
\proscons{
\begin{itemize}
  \item Crystal structure and phase identification.
  \item Lattice parameter $\to$ alloy/solid-solution composition.
  \item Crystallite size (Scherrer); bulk-averaged, representative.
  \item Simple, cheap, non-destructive.
\end{itemize}}{
\begin{itemize}
  \item Only \textbf{crystalline} phases; amorphous and very small ($<$2--3 nm) particles give no/very broad peaks.
  \item Low sensitivity to minor phases ($<$\,few \%).
  \item No information on shape, location or surface.
\end{itemize}}
